% -*- coding: utf-8 -*-


\begin{zhaiyao}
\begin{spacing}{1.5}
{
本文以沪深 300 指数增强为应用场景，对其核心量化算法——带约束的组合优化——展开优化设计与实现。指数增强的本质，是在跟踪误差、行业、个股权重与换手等约束界定的可行域内，将含噪的截面 alpha 预测稳健地转化为可实施的主动权重，即最大化主动管理基本定律 $IR = TC\cdot IC\cdot\sqrt{N}$ 中的传输系数 $TC$。本文工作分为两部分：其一，面向 A 股设计约束型多因子组合优化算法，以 LightGBM 输出 alpha、以 Ledoit-Wolf 缩减与最小特征值修正抑制协方差病态所致的``误差最大化''，并以多级求解器回退保证数值稳健；其二，针对点预测忽视个股可信度差异的缺陷，提出不确定性感知扩展，由 Conformal Prediction 导出置信度，设计不确定性进风险项与信度门控稳健混合两种稳健方案。在严格的点位成分与去前瞻口径下，长周期（2015--2024）实证表明该算法将静态多因子基线的信息比率由 $-0.99$ 提升至 $+0.20$、最大回撤由 $-77.1\%$ 收窄至 $-39.4\%$；短样本（2024--2025）中名义 $90\%$ 的 Conformal 区间实测覆盖仅约 $63\%$ 且置信度与收益倒挂，故不宜按置信度逐股加权，而信度门控稳健混合将不加权优化器的信息比率由 $-0.15$ 稳定提升至 $+0.50$（5 个随机种子全部由负转正）。本文代码与数据已开源：\url{https://github.com/aokimi0/IndexEnhancementStrategy}。
}

\end{spacing}
\end{zhaiyao}




\begin{guanjianci}
指数增强；组合优化；二次规划；传输系数；Conformal Prediction；不确定性量化
\end{guanjianci}



\begin{abstract}
\begin{spacing}{1.5}
This thesis presents the optimized design and implementation of the core algorithm behind index enhancement---constrained portfolio optimization---on the CSI~300. In essence, index enhancement transforms noisy cross-sectional alpha predictions into implementable active weights within a feasible region set by tracking-error, industry, weight, and turnover constraints, i.e., maximizing the transfer coefficient $TC$ in the fundamental law $IR = TC\cdot IC\cdot\sqrt{N}$. The work has two parts. First, a constrained multi-factor optimization algorithm for A-shares: LightGBM alpha signals, Ledoit-Wolf shrinkage with a minimum-eigenvalue correction to suppress the ``error maximization'' from ill-conditioned covariance, and a multi-solver fallback for numerical robustness. Second, an uncertainty-aware extension that remedies point prediction's neglect of per-stock confidence: Conformal Prediction supplies confidences, from which we design an uncertainty-risk term and a trust-gated robust blending. Under a strict point-in-time, look-ahead-free protocol, long-horizon (2015--2024) experiments lift the information ratio of the static multi-factor baseline from $-0.99$ to $+0.20$ and narrow the maximum drawdown from $-77.1\%$ to $-39.4\%$. On the 2024--2025 short sample, the nominal $90\%$ Conformal interval attains only about $63\%$ coverage with an inverse confidence-return relation, so per-stock weighting is unwarranted; the trust-gated robust blending instead lifts the no-weighting optimizer's information ratio from $-0.15$ to $+0.50$ (positive across all 5 seeds). Code and data are open-sourced at \url{https://github.com/aokimi0/IndexEnhancementStrategy}.
\end{spacing}
\end{abstract}


\begin{keywords}
Index Enhancement; Portfolio Optimization; Quadratic Programming; Transfer Coefficient; Conformal Prediction; Uncertainty Quantification
\end{keywords}
